\begingroup
\setlength{\emergencystretch}{2em}
\capituloTesis{Discusión}

En relación con el objetivo general, orientado a determinar la asociación entre la condición del drenaje vial y el desempeño funcional del pavimento en veinte unidades estructura--sector del corredor EMP.LI105 Pedregal--Censo--Cumbray, los coeficientes convencionales de Spearman evidenciaron diferencias entre dimensiones: CL--PCI y CL--IRI registraron 0,5940 y $-$0,6105, respectivamente, frente a 0,1714 para CT--PCI y $-$0,1053 para CT--IRI; después del ajuste de tendencia espacial y la corrección familiar de Holm, únicamente CL--PCI conservó respaldo inferencial ($\rho$ ajustado = 0,5773; p espacial = 0,0045; p Holm = 0,0180), resultado que delimita la evidencia a una relación específica dentro de las cuatro combinaciones examinadas. Esa heterogeneidad es coherente con Chu et al. (2023), quienes revisaron la interacción entre humedad, drenaje y respuesta del pavimento, así como con Winurseto et al. (2023), cuyo modelo incorporó el drenaje entre varios componentes del desempeño vial, aunque difiere en magnitud del patrón predominantemente débil comunicado por Diniz y Almeida de Melo (2023) para el drenaje superficial y el PCI urbanos; esta diferencia exige distinguir la calificación visual y los coeficientes de Pearson de aquel antecedente de los indicadores geométricos, las correlaciones por rangos y los contrastes espaciales empleados en Simbal. En conjunto, los antecedentes permiten interpretar la pertinencia de evaluar drenaje y pavimento de manera relacionada, conservando el carácter asociativo del diseño y la distinción entre el respaldo de H1 y los resultados de las demás hipótesis.

Respecto del primer objetivo específico, referido a la relación entre la evacuación longitudinal y la condición superficial del pavimento, la correlación positiva CL--PCI ($\rho$ de Spearman = 0,5940) indicó que los sectores con mayor proporción de recorrido longitudinal continuo tendieron a presentar valores superiores de PCI, relación que se mantuvo al retirar la tendencia espacial lineal ($\rho$ = 0,5773) y cuya probabilidad bilateral bajo las 9 999 aleatorizaciones de Moran fue 0,0045, con un valor ajustado por Holm de 0,0180; adicionalmente, el diagnóstico de exclusión sucesiva de una unidad conservó coeficientes positivos entre 0,5316 y 0,6807, lo que muestra que el signo de la asociación no dependió de una sola observación. El hallazgo difiere en magnitud de los coeficientes débiles publicados por Diniz y Almeida de Melo (2023), quienes obtuvieron correlaciones de Pearson de 0,39 entre pavimento y cunetas y de 0,14 entre pavimento y sumideros en diecinueve tramos de João Pessoa, aunque ambos estudios coinciden en identificar asociaciones positivas entre alguna condición de drenaje superficial y el estado del pavimento; la discrepancia de magnitudes no constituye una contradicción directa porque el antecedente evaluó subjetivamente la condición de los dispositivos urbanos y el análisis de Simbal cuantificó la continuidad de recorridos hasta su descarga; esta lectura también resulta compatible, en el plano físico y no como contraste correlacional equivalente, con la revisión de Chu et al. (2023) sobre la importancia del manejo del agua en la respuesta de las capas del pavimento y con el diagnóstico nacional de Mori Linares (2025), que documentó obstrucciones y ausencia de cunetas en la carretera Chamaya--Puente Integración. Por consiguiente, el primer objetivo específico encontró una asociación positiva respaldada por el procedimiento estadístico empleado, aunque la coincidencia entre continuidad y mejor condición superficial no demuestra que una haya producido la otra ni sustituye la evaluación de otras características constructivas del sector.

En cuanto al segundo objetivo específico, centrado en la relación entre evacuación longitudinal y rugosidad, el coeficiente convencional CL--IRI fue inverso y de magnitud apreciable ($\rho$ de Spearman = $-$0,6105; p = 0,0042), tendencia que se conservó al ajustar por coordenadas locales ($\rho$ = $-$0,5893), aunque el contraste espacial produjo p = 0,0356 y la corrección de Holm elevó ese valor a 0,1068, por encima de $\alpha$ = 0,05; en consecuencia, la hipótesis H2 no alcanzó respaldo inferencial bajo el control conjunto de las cuatro pruebas. La dirección observada resulta compatible con Winurseto et al. (2023), quienes incorporaron las condiciones del drenaje al modelamiento del desempeño vial representado por IRI y atribuyeron una contribución relativa de 21,7 \% al drenaje, porcentaje que representa una medida propia de ese modelo y, por tanto, es distinto del coeficiente de correlación estimado en Simbal; asimismo, Vega Suaña (2022) documentó la utilidad del IRI para representar el comportamiento funcional de la carretera Puno--Juliaca, con correlaciones de 0,8449 y 0,9764 correspondientes al ajuste entre predicciones y observaciones, cuya naturaleza impide equiparar esas magnitudes con una asociación drenaje--rugosidad. De esta comparación se desprende que los fundamentos técnicos de los antecedentes son compatibles con la dirección del resultado observado, mientras que la evidencia inferencial de CL--IRI permanece delimitada por el criterio de multiplicidad adoptado.

Para el tercer objetivo específico, relacionado con la abertura transversal libre y la condición superficial del pavimento, la asociación convencional CT--PCI fue directa y reducida ($\rho$ de Spearman = 0,1714; p = 0,4699), mientras que el coeficiente residual ajustado fue 0,2893, con p espacial = 0,2675 y p de Holm = 0,5350; estos valores describen una relación de baja magnitud que no obtuvo respaldo inferencial en las veinte unidades analizadas. El resultado se aproxima a las correlaciones débiles comunicadas por Diniz y Almeida de Melo (2023), particularmente a 0,14 para sumideros, aunque su estudio utilizó elementos de drenaje urbano distintos de las aberturas de entrada y salida de Simbal; por su parte, Farida et al. (2025) identificaron restricciones de capacidad hidráulica bajo determinados periodos de retorno y un PCI medio de 58,66 en Garut, en tanto que De Jager (2024) evaluó experimentalmente la modificación de caudales mediante geometrías de entrada, antecedentes que permiten contextualizar restricciones hidráulicas y condiciones superficiales, pero que no presentan una correlación directamente equivalente a CT--PCI. La comparación sitúa la falta de significancia de H3 dentro de una diferencia de indicadores y procedimientos, pues el porcentaje geométrico de abertura libre, la capacidad hidráulica y el estado superficial del pavimento corresponden a magnitudes distintas, cuya correspondencia requiere evaluaciones específicamente diseñadas para cada relación.

En el cuarto objetivo específico, referido a la abertura transversal libre y la rugosidad longitudinal, la correlación de Spearman CT--IRI alcanzó $-$0,1053 (p = 0,6587), mientras que el coeficiente residual ajustado por la ubicación fue $-$0,1400 y el contraste espacial obtuvo p = 0,5217, valor que se mantuvo sin respaldo estadístico después de la corrección de Holm (p ajustado = 0,5350); a diferencia del resultado CL--IRI, cuya magnitud fue mayor aunque tampoco superó la corrección familiar, la asociación transversal mostró una dirección inversa de escasa magnitud y el análisis de exclusión sucesiva de unidades presentó coeficientes entre $-$0,2211 y 0,0439, con cambio de signo, por lo que ni la dirección ni la significancia alcanzaron estabilidad suficiente para respaldar H4. Esta observación difiere de la relevancia atribuida al drenaje en el modelo de desempeño vial de Winurseto et al. (2023), pero no constituye una contradicción metodológicamente directa, dado que aquellos autores modelaron conjuntamente pavimento, bermas y drenaje mediante ecuaciones estructurales y utilizaron características de la sección, condición estructural y pendiente, en tanto el análisis de Simbal consideró únicamente la menor proporción de área libre entre entrada y salida de cada estructura transversal; el estudio nacional de Salas Cruz (2025) aporta una distinción complementaria al mostrar, mediante la evaluación de la carretera Canta--Huayllay, que el IRI, el PCI y la deflectometría describen propiedades diferentes de la condición vial, por lo que una característica local de la abertura de un pase de agua no necesariamente se refleja en la rugosidad agregada del sector; además, De Jager (2024) demostró experimentalmente que cambios geométricos en la entrada de una alcantarilla alteran su respuesta hidráulica bajo condiciones específicas, pero no estableció una relación con el perfil longitudinal de rodadura. Así, el cuarto resultado delimita una ausencia de evidencia estadística de asociación CT--IRI en las unidades analizadas, sin demostrar inexistencia física de interacción entre el drenaje transversal y el comportamiento del pavimento.

\endgroup
